% Seojin's report sections. Replace the corresponding sections in the final report.
\subsection{Research Framework}

We study how question selection affects population preference estimation under
an equal response budget. Human stated-preference data calibrate the MNL and
panel MMNL response generators. These generators supply the existing choice
response bank for Nemotron profiles. A survey method receives the respondent's
public profile and previously revealed answers, selects an unasked task, and
receives only that task's stored answer. Repeating this process yields nine
observations per respondent, including three prescribed initial questions.

The comparison considers fixed efficient design, Bayesian adaptive design,
non-graph adaptive selection, and graph-based adaptive selection. Respondents,
context assignments, initial questions, response realizations and question
budgets are shared across methods. The completed collections described below
cover Fixed and Non-graph. The final analysis fits a common MNL or panel MMNL
specification to the collected responses and compares estimable coefficients
with the corresponding generator parameters. Repeated appearances of a UUID
across sampling replicates are retained as part of this fixed-bank resampling
design rather than counted as distinct independent respondents.

\subsection{Question-Selection Methods}
\subsubsection{Fixed Efficient Design}

Fixed uses the original nine-task efficient-design block: Block 1 for Transfer,
Block 3 for Layover, and Block 5 for Stopover. Within each block, the prescribed
three initial tasks are presented first, followed by the remaining six tasks
in the supplied order. The respondent's answers do not change this sequence.
This baseline retains the original block's question set.

\subsubsection{Bayesian Adaptive Design}

The Bayesian comparator maintains a preference distribution for each
respondent, updates it using revealed choices, and selects subsequent tasks
using its information-based design criterion. The team's proposed
implementation uses weighted Monte Carlo particles with a Gaussian prior.
Its population prior, posterior approximation and acquisition settings must be
reported with that implementation.

\subsubsection{Non-graph Adaptive Question Selection}

The non-graph learner predicts alternative choices from respondent attributes,
the current survey context, the attributes and revealed answers of previous
questions, and the attributes of each candidate choice task. A shared event MLP
encodes each observed question--answer pair. Mean pooling produces an
order-invariant history representation, which is combined with nine respondent
covariates, the context, and history-length indicators. A shared utility head
scores alternatives A, B, C, and NONE, followed by a softmax. Neither respondent
identifiers nor question identifiers enter the network as features. We average
the predicted probabilities of five respondent-bootstrap models.

\paragraph{Offline training.}
We exclude the union of respondents in all ten evaluation replicates from both
training and validation. The learner uses 1,200 additional respondents for
training and 300 for validation, with all contexts of a respondent kept in the
same partition. Training episodes contain between zero and eight observed
answers; histories of length at least three include the prescribed initial
questions. The remaining history questions are sampled without replacement.
We minimize cross-entropy on answers outside the observed history and select
checkpoints using validation log-loss. Each ensemble member is trained for at
most 120 epochs with 1,024 episodes per epoch and early-stopping patience 30.
MNL and panel MMNL responses are used to train separate ensembles. Network
weights remain fixed throughout the evaluation surveys.

\paragraph{Acquisition rule.}
Let $h$ denote the revealed history and $\mathcal C$ the remaining candidate
questions. For a candidate $q$, we consider its four possible answers $a$ and
their current predicted probabilities $w_a(q)=p_\theta(a\mid q,h)$.
For another candidate $r$, let $p_r^{q,a}$ be its predicted choice distribution
after appending the hypothetical observation $(q,a)$ to $h$. We select the
question maximizing
\begin{equation}
\begin{aligned}
 J(q) &= \frac{1}{|\mathcal C|-1}
     \sum_{r\in\mathcal C\setminus\{q\}}\left[
       H\!\left(\sum_a w_a(q)p_r^{q,a}\right)\right.\\[-2pt]
 &\hspace{35mm}\left.-\sum_a w_a(q)H(p_r^{q,a})\right].
\end{aligned}
\end{equation}
This weighted Jensen--Shannon score measures changes in predictions for other
remaining questions. Its value is used as a predictive-dependence proxy;
coefficient recovery is assessed separately by the common final estimator.
Hypothetical answers never modify the actual respondent history.

\paragraph{Survey collection.}
Each replicate contains 500 respondents in each of the three contexts. All
methods use the supplied context assignments and the same three initial
questions and answers. Non-graph selection then asks six additional questions
from the fifteen remaining questions in the same context. Fixed selection uses
the supplied order of the original nine-question block. Answers are read from
the existing response bank only after a question is selected. The ten
replicates resample respondents from this fixed bank. All nine revealed answers,
including the initial three, are retained for the common MNL or panel MMNL
coefficient-recovery analysis. No prediction holdout is used in this experiment.

\subsubsection{Graph-Based Adaptive Selection}

The graph comparator represents respondents and question--alternative pairs
as nodes connected by observed response relations. Message passing shares
response information across connected nodes before predicting choices.
The documented comparison uses the same weighted JS acquisition strategy as
Non-graph. A controlled comparison matches the available profile and task
information, revealed-history restrictions, choice-prediction objective and
data partitions. The team's preliminary graph study selected LightGCN and
examined different propagation depths; final architecture and masking details
belong to the graph implementation description.

\subsection{Completed Survey Collections}

For each of ten supplied replicates, Fixed and Non-graph collected nine
answers from 500 respondents in each of three contexts under each generator.
This produced 40 response files and 540,000 respondent--task observations.
All files passed the supplied schema and survey validator. Additional checks
matched the original cohort fields and stored response-bank answers,
verified the prescribed initial sequence and context quotas, and confirmed
that no evaluation UUID entered Non-graph training or validation.

These are the collected response sets for the shared final estimator.
Coefficient-recovery results will report bias and RMSE for the estimable mean
coefficients and, for panel MMNL, the random-coefficient standard deviations.
The population estimator and normalization must be common across methods.
